% Articulación posterior a la generación de pi, phi y e.
% Contenido acumulativo; no reemplaza el desarrollo previo del TRIT.
\subsection{Réalisations exponentielles des coordonnées générées}
\label{euler-trit-articulacion}

Les coordonnées régionales déjà déterminées permettent de reconnaître les
réalisations concrètes de la collection quadratique du TRIT. L'ordre est
essentiel : la construction des régions précède l'évaluation d'une
exponentielle aux paramètres \(\pi\) ou \(2\log\varphi\). Ces expressions
identifient la fonction des coordonnées dans des opérateurs construits à partir
d'APP et du TPK ; elles n'interviennent pas dans la sélection des préfixes qui les produisent.

\subsubsection{Cycle ternaire, transport nilpotent et incidence des feuillets}

La multiplication \(a(r)=4r\) dans \(\mathbb Z/9\mathbb Z\) est d’ordre
\(3\). Sur les unités, elle détermine les orbites \(\{1,4,7\}\) et
\(\{2,8,5\}\). Dans le module d’augmentation de l’une quelconque d’entre elles, formé
des combinaisons entières dont les coefficients ont une somme nulle, la base
\((e_r-e_{4r},e_{4r}-e_{7r})\) donne
\[
 C=\begin{pmatrix}0&-1\\1&-1\end{pmatrix},
 \qquad C^2+C+I=0.
\]
La forme restreinte a pour matrice de Gram
\(\bigl(\begin{smallmatrix}2&-1\\-1&2\end{smallmatrix}\bigr)\).
Dans les évaluations résiduelles modulo \(9\), l’action distingue les feuillets :
\(S(ax,ay)=aS(x,y)\), tandis que \(P(ax,ay)=a^{-1}P(x,y)\).
La réalisation cyclique conserve donc une orientation contragrédiente
entre somme et produit. Les quotients du relèvement entier sont conservés
dans le registre arithmétique correspondant.

Le transport parabolique élémentaire est représenté par
\[
 N=\begin{pmatrix}0&1\\0&0\end{pmatrix}.
\]
L’incidence surfacique--volumique déjà construite dans \eqref{eq:fav} fournit
\[
 F_{\rm av}=\begin{pmatrix}0&1\\1&1\end{pmatrix},
 \qquad A=F_{\rm av}^2=\begin{pmatrix}1&1\\1&2\end{pmatrix}.
\]
Les trois opérateurs agissent dans des espaces réels déclarés : le module
d’augmentation tensorisé avec \(\mathbb R\), le plan du transport nilpotent et
le plan d’incidence modale. Le partage d’une dimension matricielle n’identifie
ni ces espaces ni leurs coordonnées.

\begin{proposition}[Générateurs concrets des trois types]
\label{euler-trit-generadores}
Les opérateurs
\[
 J_{\mathrm e}=\frac{2C+I}{\sqrt3},\qquad
 J_{\mathrm p}=N,\qquad
 J_{\mathrm h}=\frac{2A-3I}{\sqrt5}
\]
satisfont, respectivement,
\[
 J_{\mathrm e}^{\,2}=-I,\qquad
 J_{\mathrm p}^{\,2}=0,\qquad
 J_{\mathrm h}^{\,2}=I.
\]
\end{proposition}
\begin{proof}
La relation de \(C\) implique \((2C+I)^2=-3I\).
La multiplication directe donne \(N^2=0\).
Enfin, \(A\) est de trace \(3\) et de déterminant \(1\), d’où
\(A^2-3A+I=0\) et \((2A-3I)^2=5I\).
\end{proof}

\begin{proposition}[Trois évaluations de l’identité exponentielle]
\label{euler-trit-evaluaciones}
Dans les réalisations précédentes,
\begin{equation}
 \boxed{
 \exp(\pi J_{\mathrm e})+I=0,\qquad
 \exp(J_{\mathrm p})=I+J_{\mathrm p},\qquad
 \exp(2\log\varphi\,J_{\mathrm h})=F_{\rm av}^{\,2}.}
 \label{euler-trit-tres-evaluaciones}
\end{equation}
\end{proposition}
\begin{proof}
La spécialisation elliptique de l’exponentielle tritique donne
\(\cos\pi\,I+\sin\pi\,J_{\mathrm e}=-I\).
La nilpotence élimine exactement les termes de degré au moins \(2\)
dans la deuxième expression. Pour la troisième, \(\varphi^2=\varphi+1\) implique
\[
 \cosh(2\log\varphi)=\frac32,\qquad
 \sinh(2\log\varphi)=\frac{\sqrt5}{2}.
\]
L’exponentielle vaut donc
\(\frac32I+\frac{\sqrt5}{2}J_{\mathrm h}=A\).
\end{proof}

La première égalité représente le demi-tour elliptique ; le tour complet
satisfait \(\exp(2\pi J_{\mathrm e})=I\). La deuxième maintient un terme
linéaire non nul : le régime parabolique exprime un transport, et non une absence
d'évolution. La troisième réalise une avancée hyperbolique unimodulaire. Ses
valeurs propres sont \(\varphi^2\) et \(\varphi^{-2}\) ; une direction se dilate et
l'autre se contracte en conservant le déterminant. Les trois évaluations
emploient des paramètres distincts d'une collection exponentielle commune.

Dans cette collection, \(e\) intervient par l'évaluation scalaire
\(\exp(I)=eI\) et la loi \(\exp(sI)\exp(tI)=\exp((s+t)I)\).
Son rôle ne se limite pas au type parabolique. La coordonnée \(e\), obtenue
auparavant à partir de sa région, est reconnue ici comme l'évaluation en l'unité
de l'exponentielle ; cette identité ne remplace pas sa généalogie coinductive.

\subsubsection{Résolution polynomiale des régimes}

Les trois réalisations admettent une présentation compacte sur la somme
directe de leurs espaces. On définit
\[
 \mathbb J=J_{\mathrm e}\oplus J_{\mathrm p}\oplus J_{\mathrm h},
 \qquad Q=\mathbb J^2.
\]
Les relations précédentes impliquent
\[
 \mathbb J^6=\mathbb J^2,\qquad Q^3=Q.
\]
L’emploi de \(Q\) maintient ce carré opératoriel distinct du registre
dodécaphasique \(K\).

\begin{proposition}[Projecteurs polynomiaux de régime]
\label{euler-trit-proyectores}
Les applications
\[
 p_{\mathrm e}=\frac{Q^2-Q}{2},\qquad
 p_{\mathrm p}=I-Q^2,\qquad
 p_{\mathrm h}=\frac{Q^2+Q}{2}
\]
sont des idempotents mutuellement orthogonaux, de somme \(I\), qui restituent les trois
termes de la somme directe représentant les régimes.
\end{proposition}
\begin{proof}
Les trois polynômes prennent les valeurs \((1,0,0)\), \((0,1,0)\) et
\((0,0,1)\) aux points \(-1,0,+1\), respectivement. Les différences
entre leurs carrés et eux-mêmes, ainsi que leurs produits croisés, sont
des multiples de \(X^3-X\). Substituer \(Q\), qui annule ce polynôme, démontre
les identités. Dans les trois blocs, \(Q\) agit comme \(-I,0,I\).
\end{proof}

On obtient ainsi
\begin{align}
 \exp(t\mathbb J)
 ={}&p_{\mathrm e}(\cos t\,I+\sin t\,\mathbb J)
   +p_{\mathrm p}(I+t\mathbb J)\notag\\
   &+p_{\mathrm h}(\cosh t\,I+\sinh t\,\mathbb J).
 \label{euler-trit-exponencial-total}
\end{align}
Chaque terme conserve sa relation quadratique et la conjugaison inverse
son évolution. La représentation totale réunit les trois régimes, mais ne
définit pas à elle seule un opérateur de transition entre eux. Une telle transition
requiert les applications de transport du TPK avec leurs domaines, leur orientation et
leur mémoire. L’articulation exponentielle permet de reconnaître leurs types sans
remplacer cette dynamique par une identification des fibres.

% PROCEDENCIA FOCAL (no impresa)
% Resultado recuperado: tres generadores y tres evaluaciones exponenciales.
% Propietario completo:
% BASE/manuscrito/sections/hmt/09b_algebras_cuadraticas.tex:289-516.
% Antecedente causal de C y acción contragrediente:
% BASE/manuscrito/sections/hmt/03d_esqueleto_ciclotomico_app.tex:13-47,145-246.
% Los tres tipos y la orientación están en:
% BASE/manuscrito/sections/hmt/02_trit_geometrias.tex:83-196.
% Resolución polinómica recuperada:
% output/INVESTIGACION_HOLONOMIA_NONADICA_CRISTAL_APERIODICO_20260903/
% ALGEBRA_HOLONOMICA_UNIVERSAL_Y_EULER_TRITICO.md, secciones 8-11 y 22.
% Los cuatro propietarios se leyeron completos. BASE designa:
% /Users/ruben/Documents/HMT2/
% EL_CIERRE_HOLOGRAFICO_DEL_INFINITO_SUCESOR_102_CAPITULOS_2026-08-28_EN_TRABAJO/
% No se incorpora la identidad P9(alpha)=delta del documento especializado:
% esa identificación no pertenece a este bloque.
% Tampoco se identifica una dirección nilpotente con la vacancia electrónica,
% ni se presume una representación global del transporte por la suma directa.
% COMPROBACIONES: Python estándar, enteros y Fraction.
% C^2+C+I=0; (2C+I)^2=-3I; N^2=0; (2Fav^2-3I)^2=5I.
% 81 pares de residuos verifican las dos acciones contragredientes.
% Los 3 proyectores se verificaron en Q=-1,0,+1; prueba general en texto.
% 6 labels únicos y entornos equilibrados. Sin compilación.
